\documentclass[12pt]{article}
\usepackage[utf8]{inputenc}
\usepackage[T1]{fontenc}
\usepackage{hyperref}
\title{Fagani Language Resources for K-12 Curriculum Development}
\author{Compiled by Claude Sonnet 5.5 (claude-sonnet-5-5) for Andr\'as Kornai}
\date{}
\begin{document}
\maketitle

\section{Introduction}
Fagani (ISO 639-3: faf; Glottocode faga1239) is an Austronesian language of northwest Makira (San Cristobal), Solomon Islands, in the Southeast Solomonic branch. The usual figure is about 900 speakers (1999); Joshua Project lists 1,900 ethnic Fagani. Dialects are Agufi, Fagani and Rihu'a. UNESCO classes it as vulnerable. Little published material was found.

\section{Primary Educational Resources}
No primers, textbooks or school programmes in Fagani were found in search results.

\section{Linguistic and Reference Materials}
\begin{itemize}
\item Wikipedia, Fagani language. \url{https://en.wikipedia.org/wiki/Fagani_language}
\item Verbix page for Fagani. \url{https://docs.verbix.com/Languages/Fagani}
\end{itemize}

\section{Religious and Community Resources}
\begin{itemize}
\item Joshua Project profile of the Fagani people; it reports Bible portions exist but not a full New Testament. \url{https://joshuaproject.net/people_groups/11176}
\end{itemize}

\section{Digital and Audio Resources}
The search summary mentioned materials in the Arthur Capell collection held by Paradisec, but no URL for it appeared in the results, so none is cited.

\section{Curriculum Development Context}
The language is vulnerable and not known to be taught in schools. Curriculum work would require new documentation and community-led materials; recent speaker surveys were not found.

\end{document}